\chapter{Preliminares}
\label{ch:preliminares}

Este capítulo reúne el contexto en la literatura y los fundamentos teóricos sobre
los que se apoya el trabajo. Su propósito es servir de referencia: los capítulos
de desarrollo emplean estas herramientas en su forma aplicada y remiten aquí, en
lugar de reexplicarlas. Se presenta cada familia de modelos junto con el papel que
desempeña en el problema, y se reserva para los capítulos correspondientes la
justificación de las decisiones concretas (tamaño de celda, número de estados,
hiperparámetros) y la definición de las métricas de evaluación.

\section{El análisis y la predicción del movimiento animal}
\label{sec:prelim-movimiento}

El movimiento es uno de los procesos que articulan la vida de un animal, pues
conecta su fisiología y su comportamiento con el entorno en el que habita. La
ecología del movimiento propone un marco unificador para estudiarlo, que sitúa el
desplazamiento de un organismo en el centro y lo relaciona con su estado interno,
sus capacidades de navegación y locomoción y los factores externos que lo
condicionan \cite{nathan2008}. La generalización de los dispositivos de
geolocalización ha convertido esta disciplina en un campo intensivo en datos:
repositorios como Movebank acumulan grandes volúmenes de trayectorias y han
abierto el análisis del movimiento a las técnicas estadísticas y computacionales
\cite{kays2015, demsar2015}.

Con esa abundancia de datos, el interés se desplaza de describir trayectorias
pasadas a modelarlas y anticipar posiciones futuras. Esa predicción exige combinar
herramientas de naturaleza distinta, en lugar de depender de un único modelo, por
la dificultad ya señalada del problema. Las secciones siguientes presentan los
fundamentos de las que emplea este trabajo.

\section{Cadenas de Markov}
\label{sec:prelim-markov}

Una cadena de Markov es el modelo probabilístico más sencillo para una secuencia de
estados en la que el futuro depende del pasado solo a través del presente. Aplicada
al movimiento, describe el desplazamiento entre zonas discretas del espacio con
pocos parámetros y una interpretación directa, lo que la convierte en una
referencia natural frente a la que medir modelos más complejos.

Formalmente, sea $\{X_t\}_{t \geq 0}$ un proceso estocástico que toma valores en un
espacio de estados finito $S = \{1, \dots, N\}$. El proceso es una cadena de Markov
de primer orden si satisface la propiedad de Markov, es decir, si el estado
siguiente depende únicamente del estado actual y no de toda la historia previa:

\begin{equation}
  P(X_{t+1} = j \mid X_t = i,\, X_{t-1} = i_{t-1}, \dots, X_0 = i_0)
  = P(X_{t+1} = j \mid X_t = i),
  \label{eq:prelim-markov}
\end{equation}

para todo par de estados $i, j \in S$. Cuando esta probabilidad no depende del
instante $t$, la cadena es homogénea y queda completamente determinada por la
matriz de transición $P = (p_{ij})$, con $p_{ij} = P(X_{t+1} = j \mid X_t = i)$.
Cada fila de $P$ es una distribución de probabilidad sobre los estados de destino,
de modo que $p_{ij} \geq 0$ y $\sum_{j} p_{ij} = 1$ para todo estado de origen
\cite{norris1997}. La probabilidad de transitar entre dos estados en $n$ pasos es
la entrada correspondiente de la potencia $P^{n}$ (ecuaciones de
Chapman-Kolmogorov), lo que permite propagar la distribución de estados a varios
pasos vista.

Las probabilidades de transición no se conocen de antemano, sino que se estiman a
partir de las transiciones observadas. El estimador de máxima verosimilitud de
$p_{ij}$ es la frecuencia relativa de las transiciones del estado $i$ al $j$,

\begin{equation}
  \hat{p}_{ij} = \frac{n_{ij}}{\sum_{k} n_{ik}},
  \label{eq:prelim-mle}
\end{equation}

donde $n_{ij}$ es el número de transiciones observadas de $i$ a $j$. El capítulo
\ref{ch:o2} parte de este estimador y lo ajusta para que los destinos aún no
observados no reciban probabilidad nula.
